\section{Purpose}
\label{sec:purpose}%

Nowadays, technology has aided every aspect of our lives, and has been involved into every workspace, hobby and activity, facilitating them. \textbf{Best Bike Paths (BBP)} serves as a simple yet useful way for bikers to easily track and plan their routes. 

This document explains the design decisions made for the development of \textbf{Best Bike Paths (BBP)} and guides the aspects of the application and its system.

\subsection{Scope}
\label{subsec:scope}%

The system aims for an easy interaction for bikers when intending to search new paths to traverse, or to travel paths they already know. It allows bikers to rate paths for other bikers to see and select their state in case they are in bad condition. Users can also see their past trips accompanied with their data, and, if possible, the weather conditions during the trip. Furthermore, the system is able to detect when the user is on their bicycle even when the trip has not been manually activated, this helps users track daily and frequent paths. With these features, \textbf{Best Bike Paths (BBP)} aims to be an easy and useful application for bikers that will make their life easier.

\section{Definition, Acronyms, Abbreviations}
\label{sec:definition_acronyms_abbreviations}%

\subsection{Definitions}
\label{subsec:definitions}%
\begin{itemize}
    \item \textbf{Bike Path:} A path or marked route specifically designed or suitable for people riding bicycles.
    \item \textbf{Route:} A specific path or course taken by a cyclist from an origin point to a destination point.
    \item \textbf{Trip:} A recorded journey made by a cyclist, including start time, end time, duration, distance, and the path followed.
    \item \textbf{Obstacle:} Any physical impediment or hazard on a bike path, such as potholes, debris, or other obstructions that may affect cycling safety or comfort.
    \item \textbf{Registered User:} A user who has created an account in the system, enabling access to advanced features such as trip recording, path publishing, and personal statistics.
    \item \textbf{Non-Registered User:} A user who accesses the system without creating an account, with limited access to browse and search available paths.
    \item \textbf{Automatic Detection:} The system's capability to automatically identify when a user is cycling and to detect paths and obstacles using sensor data from mobile devices.
    \item \textbf{Path Status:} A classification indicating the condition of a bike path, typically categorized as optimal, medium, sufficient, or requiring maintenance.
    \item \textbf{Path Score:} A computed metric that ranks bike paths based on their status, effectiveness, and other relevant factors to help users choose the best route between two points.
\end{itemize}

\subsection{Acronyms}
\label{subsec:acronyms}%
\begin{itemize}
    \item \textbf{BBP:} Best Bike Paths
    \item \textbf{GPS:} Global Positioning System
    \item \textbf{UI:} User Interface
    \item \textbf{API:} Application Programming Interface
    \item \textbf{DBMS:} Database Management Systems
\end{itemize}

\subsection{Abbreviations}
\label{subsec:abbreviations}%
\begin{itemize}
    \item \textbf{UC:} Use Case
    \item \textbf{F:} Functional Requirement
\end{itemize}

\section{Revision History}
\label{sec:revision_history}

The document has undergone several revisions throughout its development:

\begin{itemize}
    \item \textbf{Version 0.1} (2 January 2025): Initial Release by José Martínez Estévez and Pablo Jesús Herrero Díaz
    \item \textbf{Version 0.2} (5 January 2025): Architectural Design, User Interface Design by José Martínez Estévez and Pablo Jesús Herrero Díaz
    \item \textbf{Version 0.3} (6 January 2025): Requirements Traceability, Implementation Integration and Test Plan by José Martínez Estévez and Pablo Jesús Herrero Díaz
    \item \textbf{Version 1.0} (7 January 2025): Final version by José Martínez Estévez and Pablo Jesús Herrero Díaz
\end{itemize}


\section{Reference Documents}
\label{sec:reference_documents}%

\begin{itemize}
    \item Specification for the 2025/2026 Design Document
    \item Course slides available on the WeBeep page 
    \item Previous year student project available on the WeBeep page
\end{itemize}

\section{Document Structure}
\label{sec:document_structure}%
The document is organized into the following sections, each serving a distinct purpose:

\paragraph{1. Introduction:} This section presents the project's purpose, scope, and essential information required for understanding the project domain correctly. It includes definitions, acronyms, and abbreviations to facilitate comprehension of the problem domain.

\paragraph{2. Architectural Design:} Delivers a detailed description of the design of the application system.

\paragraph{3. User Interface Design:}  This section indicates the planning for the UI.

\paragraph{4. Requirements Traceability:} Explains how the goals defined in the RASD are covered by the features offered in the application.

\paragraph{5. Implementation, Integration and Test Plan:} This section covers how the application must be implemented and some guidelines for the testing of it.

\paragraph{6. Effort Spent:} Describes the contribution of each team member to the writing of this document.

\paragraph{7. References:} Provides a list of all supplementary materials and references utilized in the production of this document.

